% The editable source of report.pdf. Build with:  tectonic report.tex
% Every number in this document is produced by a stage under ../stages/ and is
% written to ../outputs/tables/ and ../outputs/diagnostics/ by that stage.
\documentclass[11pt,a4paper]{article}
\usepackage[margin=2.4cm,top=2.6cm,bottom=2.6cm]{geometry}
\usepackage{booktabs}
\usepackage{graphicx}
\usepackage{amsmath,amssymb}
\usepackage{microtype}
\usepackage{caption}
\usepackage{array}
\usepackage{longtable}
\usepackage[table]{xcolor}
\usepackage{titlesec}
\usepackage[colorlinks=true,linkcolor=darkblue,citecolor=darkblue,urlcolor=darkblue]{hyperref}
\definecolor{darkblue}{HTML}{2B6CA3}
\definecolor{warm}{HTML}{C1553B}
\definecolor{shade}{HTML}{F4F4F2}
\captionsetup{font=small,labelfont=bf,width=0.95\textwidth}
\titleformat{\section}{\Large\bfseries}{\thesection}{0.7em}{}
\titleformat{\subsection}{\large\bfseries}{\thesubsection}{0.6em}{}
\setlength{\parskip}{0.45em}
\setlength{\parindent}{0pt}
\newcommand{\good}[1]{\textcolor{darkblue}{\textbf{#1}}}
\newcommand{\bad}[1]{\textcolor{warm}{\textbf{#1}}}
\newcommand{\figref}[2]{\begin{figure}[htbp]\centering\includegraphics[width=\textwidth]{../outputs/figures/#1.pdf}\caption{#2}\end{figure}}

\title{\vspace{-1.2cm}\textbf{The stochastic structure of Hong Kong racing performance}\\[0.35em]
\large An exploratory latent-performance study, and what of it belongs in a betting model}
\author{Alternative Performance-Distribution Research Project}
\date{31 August 2026}

\begin{document}
\maketitle
\vspace{-0.8cm}

\begin{center}
\fcolorbox{black!25}{shade}{\parbox{0.94\textwidth}{\small
\textbf{What this document is.} An exploratory study of the distribution of horse
performance in HKJC Class 1--5 racing, 2008--2026, built from individual finishing
times and per-section sectional times. It is a companion to an existing production
winner-probability model, which it treats as read-only and as the reference for
what a useful result would have to look like.

\textbf{What it does not establish.} Nothing here is a profitability claim. Every
market price used is a result-page price, recorded after the pool closed, at which
no bet could have been placed; no dividends were loaded and no return is computed
anywhere in this project. Every season in the archive has been development data,
here and in the reference project, so no window is genuinely unseen. The
confirmation step this study cannot substitute for is a prospective test of a
locked model against timestamped pre-race price snapshots.}}
\end{center}

\tableofcontents
\newpage

%======================================================================
\section{Executive summary}
%======================================================================

The reference production model beats the market by \good{$-0.0045$} of race log
loss. That number is the yardstick for everything below: a piece of performance
structure is worth having only if it moves a meaningful fraction of it.

Fourteen structural questions were asked. The answers divide unevenly, and the
division is the result.

\subsection*{What is real, and useful}

\begin{itemize}\itemsep2pt
\item \textbf{Sections are exactly comparable across distances.} HKJC times every
race in 400\,m sections measured back from the finish, the first carrying the
remainder --- verified on all 13{,}309 races in the population. So ``the last
400\,m'' is one physical quantity in a 1200\,m sprint and a 2200\,m staying race.
Both of the reference project's own sectional engines read the first and last
section only; the middle of the race was unexamined, and this study examines it.

\item \textbf{Jockey ability is identifiable, and is a far better forward
predictor than horse ability.} Jockey identity alone forecasts relative
performance at out-of-sample $R^2$ \good{0.048} against horse identity's 0.012,
and at a forward scale of \good{1.05} against the horse effect's 0.39. Not because
it is larger --- it is a fifth the size --- but because it is far better
determined: about 1{,}000 rides per jockey against 15 starts per horse, and no
drift. Two independent estimators agree at Spearman 0.970.

\item \textbf{Horse ability drifts with a half-life of about two months.} Forward
correlation is monotone in recency and peaks at 60 days (0.150 against a static
0.114). A static ability estimate is over-scaled by a factor of 2.7.

\item \textbf{Horse-specific volatility is real and forecastable.} 61.8\% of the
observed spread in per-horse log residual standard deviation is signal rather than
sampling noise; a horse one standard deviation more volatile has 1.43$\times$ the
spread, and the property persists out of sample at $z\approx 8$.

\item \textbf{The bad tail is heavier, and is partly a recorded event.} The
out-of-sample residual has skewness $-0.729$ and excess kurtosis $+0.810$. Of runs
more than three standard deviations below expectation, \good{46\%} carry a severe
stewards' phrase against a 9.5\% base rate --- a lift of 4.8, monotone across all
seven residual bands.

\item \textbf{Race-level shared shocks are large, and track speed is a property of
the day.} The shared shock is 32\% of absolute residual variance, about half a
second at 1200\,m. A running mean of earlier races on the same card correlates
\good{0.50} with a race's own shock, rising to 0.58 by the eighth race of a
meeting.
\end{itemize}

\subsection*{What is real, and not useful}

\begin{itemize}\itemsep2pt
\item \textbf{Horse $\times$ environment sensitivity exists only for some
variables.} Venue (Happy Valley against Sha Tin) and surface sensitivity persist
out of sample at $z=8.3$ and $z=6.4$. Going sensitivity does \bad{not}
($z=1.37$) --- even though going moves absolute race time by 1.2\,s. This is
precisely the outcome the research plan nominated in advance as its most valuable
possible finding, and it lands.

\item \textbf{The sectional profile is more persistent than final time, for the
wrong reason.} Fourteen sectional quantities beat the best final-time target on
persistence, the strongest reaching forward $r=0.483$ against $0.278$. But
persistence and informativeness are nearly \emph{inversely} ordered: the same
quantity predicts its own next value at 0.483 and next \emph{performance} at
0.045. What persists is running \emph{style}, a habit; style persists precisely
because it carries little ability information.

\item \textbf{The market underprices lightly-raced horses, and the production
model already knows.} A measure of ability uncertainty essentially uncorrelated
with the price ($-0.004$) predicts the outcome given the price at $t=+15.0$. But
the incumbent's own 28 columns already reconstruct \bad{73\%} of it. There is
nothing left to add.
\end{itemize}

\subsection*{What is not there}

\begin{itemize}\itemsep2pt
\item \textbf{Horse $\times$ jockey compatibility is not forecastable.} The median
pair has \bad{one} ride and 87.8\% have fewer than five. On pairs with enough
rides to test, the early-to-late correlation is indistinguishable from zero
($z$ from $-0.4$ to $+0.8$) on a design that would have shown $r=+0.21$ for a true
pair standard deviation of 0.20. The \emph{in-sample} fit meanwhile reports a
10.3\% variance share --- a clean illustration of why training fit is not evidence.

\item \textbf{The draw effect is large and fully priced.} Relative performance
falls 0.25 standard deviations from the inside to the outside draw; the market
residual across the same bands is non-monotone and within 0.2 percentage points of
zero, and a date-atomic draw-bias estimate predicts the market's error at
$r=0.007$.

\item \textbf{No elaboration of the shock distribution improves the
probabilities.} Once a single utility scale is fitted per variant, plain
homoskedastic normal shocks are the best of six nested generative variants ---
heteroskedasticity, horse-specific volatility, Student-$t$ tails, latent-state
uncertainty and a shared pace factor all fail to improve on it. This
independently reproduces the reference project's own probit result that the
utility scale is the whole game, now with measured heteroskedasticity and a
sectional-derived pace loading rather than a position habit.
\end{itemize}

\subsection*{The practical answer}

Nine candidate feature blocks were built forward-only and tested against the
production architecture on identical rows, folds and seeds, with a no-op control
that returns \good{exactly zero}. \textbf{None clears a promotion gate.} Three are
directionally favourable at 20--26\% of the production edge, with $|t|\approx 1$
and three of five years improving --- \emph{investigate}, not \emph{promote}, on
the reference project's own criteria.

And the pattern across the nine is the most useful single thing this study
produced: the correlation between how \emph{novel} a candidate is and how much it
\emph{helped} is $+0.39$, meaning \textbf{the more novel the information, the less
it helped}. Marginal value at this effect size lies in sharpening what the model
already uses, not in adding what it lacks. That is a coherent explanation for the
reference project's own history: roughly four hundred candidate columns built,
audited, and none promoted.

\subsection*{The two things that would change the answer}

\begin{enumerate}\itemsep2pt
\item A canonical-scale confirmation of the sectional finishing block --- five
seeds, 400 rounds, full chronology --- because the development-profile estimate
($-0.00088$, 20\% of the production edge) is favourable and its standard error is
of the same order as its mean.
\item The \textbf{within-card track-speed signal}. It is the one measurement here
that is both large ($r=0.50$) and \emph{structurally unavailable} to the reference
pipeline, which excludes all same-date information by rule. That rule exists to
stop a horse seeing its own same-day result; a track-speed reading taken from race
one before race eight runs is forward in time and is not that.
\end{enumerate}

\newpage
%======================================================================
\section{The existing project and the data}
%======================================================================

\subsection{What the reference system does}

The reference project estimates a calibrated probability that each declared runner
wins, for HKJC Class 1--5 races at Sha Tin and Happy Valley. Its production
estimator does not learn a probability; it learns a \emph{correction} to the
market's own. For race $r$ and runner $i$, with $q_{ri}$ the market's normalised
implied probability,
\[
u_{ri}=\log q_{ri}+c_\theta(x_{ri}),
\qquad
p_{ri}=\operatorname{softmax}_i(u_r),
\]
where $\log q$ enters XGBoost as a fixed additive base margin and the booster
learns $c_\theta$ against the racewise winner-softmax likelihood. At $c=0$ the
model \emph{is} the market exactly, because $q$ already sums to one within the
race. Every published improvement is therefore an improvement over the market on
the same races.

Its measured standing, on 7{,}933 races over 2016--2026: the production five-seed
ensemble beats the market by $-0.004533$ of race log loss in ten of eleven years.
The unanchored baseline beats it by $-0.000992$ but changes sign in four years.
That contrast --- anchoring converts a coin flip into a consistent one --- is the
architecture's argument.

One prior result is directly relevant here and is the strongest existing evidence
for the framing this study adopts. Replacing independent Gumbel choice shocks with
independent normal ones, holding the mean utilities fixed, improves race log loss
by $-0.0126$ standalone and $-0.00094$ market-relative, in eleven of eleven years.
\textbf{Latent performance with normal noise is already known to describe this
sport better than the softmax the production model uses.}

\subsection{The archive, and the measurement that makes this study possible}

Fourteen scraped datasets are available. This study uses five: \texttt{sectional-times}
(192\,MB) for per-section times, \texttt{aspx-results} for prices, weights, draw
and human connections, the \texttt{horses} per-start JSON for the official rating
and declared gear (which the result pages do not carry), \texttt{horse-profiles},
and \texttt{comments-on-running} for the stewards' text.

\begin{table}[htbp]\centering\small
\caption{Sectional geometry, measured on every season. Sections after the first
are exactly 400\,m, so indexing them \emph{from the finish} makes them physically
comparable across every distance. Valid on 1.0000 of races.}
\begin{tabular}{rrl}
\toprule
Distance (m) & Sections & Section lengths, start to finish\\
\midrule
1000 & 3 & 200, 400, 400\\
1200 & 3 & 400, 400, 400\\
1400 & 4 & 200, 400, 400, 400\\
1600 & 4 & 400, 400, 400, 400\\
1650 & 4 & 450, 400, 400, 400\\
1800 & 5 & 200, 400, 400, 400, 400\\
2000 & 5 & 400, 400, 400, 400, 400\\
2200 & 6 & 200, 400, 400, 400, 400, 400\\
\bottomrule
\end{tabular}
\end{table}

Coverage is close to complete: over 648{,}121 section rows,
\texttt{sectional\_time\_sec} is missing on 0.0000 in every one of nineteen
seasons, finishing time on 0.0002--0.0017, and the per-section running position on
0.0000--0.0004. Each row also carries the leader's own time for that section,
the margin behind the leader, and --- for sections after the first --- 200\,m
sub-splits, which this study has not exploited.

\subsection{Population, and reconciliation against the reference}

A race enters if it is Class 1--5 at one of the two Hong Kong tracks, every runner
is priced, the reciprocal-odds sum is at least 1.10, the field has at least four
runners and the finishing time is physically plausible. Races are kept whole or
dropped whole, because a partial field breaks every within-race normalisation.

Result: \textbf{165{,}435 runner-races over 13{,}309 races}, 2 April 2008 to
15 March 2026, mean overround 1.2291.

Restricting to 2010 onward to match the reference training start, 12{,}041 of its
12{,}237 races are shared, and on 149{,}494 jointly identified runner rows the
market probability computed here agrees with the reference table's to
\good{$5.55\times10^{-17}$} at the 99th percentile --- the same number to float64
noise.

The residual disagreement is understood rather than tolerated. Exactly 72 races
(0.60\%) have one fewer runner here, never more, and the missing runners are
\texttt{UR}, \texttt{PU} and \texttt{WV} codes: \textbf{the sectional archive
records no times for a non-finisher}. The reference project deliberately keeps
non-finishers in the pre-race field as losing starters, because removing them by
post-race status is field-definition look-ahead. This study's performance panel is
therefore conditioned on finishing, which is a selection effect carried into
\S\ref{sec:limitations}.

\newpage
%======================================================================
\section{Constructing a performance target}
%======================================================================

\subsection{Model 0: what the environment explains}

A ridge regression of performance on one-hot environment blocks, with no horse,
jockey or market column in the design --- a par that knew who was running would
absorb the ability it exists to measure. Two versions are kept apart: an in-sample
fit used only to decompose variance, and a forward fit, refitted for each
evaluation year on strictly earlier years, used for every predictive claim.

\begin{table}[htbp]\centering\small
\caption{Nested in-sample $R^2$. Almost all of raw \emph{time} is distance; the
interesting variation lives in speed.}
\begin{tabular}{lrrrr}
\toprule
Block & $R^2$ (seconds) & $\Delta R^2$ & $R^2$ (log speed) & $\Delta R^2$\\
\midrule
distance & 0.994597 & 0.994597 & 0.661373 & 0.661373\\
$+$ venue $\times$ surface & 0.994831 & 0.000234 & 0.677483 & 0.016110\\
$+$ going & 0.995291 & 0.000460 & 0.707558 & 0.030075\\
$+$ course \& rail & 0.995312 & 0.000020 & 0.708825 & 0.001267\\
$+$ class & 0.995781 & 0.000470 & 0.740827 & 0.032002\\
$+$ field size & 0.995795 & 0.000014 & 0.741852 & 0.001025\\
$+$ season & \textbf{0.996177} & 0.000382 & \textbf{0.767558} & 0.025706\\
\bottomrule
\end{tabular}
\end{table}

Forward par $R^2$ on log speed averages 0.7479 across evaluation years, so the par
model generalises about as well as it fits.

\subsection{Beaten lengths are not a measurement}

The plan lists ``beaten-length-adjusted performance'' as a candidate target. One
measurement removes it. Regressing beaten seconds on beaten lengths through the
origin gives a slope of \textbf{0.15996\,s per length}, and the same value to four
decimals on every distance from 1000\,m to 2200\,m --- that is exactly 6.25
lengths per second. Testing the stronger claim:
\[
\text{lengths} = \operatorname{round}\!\big(\text{beaten seconds}\times 6.25\times 4\big)/4
\]
holds \good{exactly} on 83.7\% of 148{,}575 rows and to within a quarter length on
99.7\%; the median ratio is 6.250000.

\texttt{LBW} is therefore the recorded time restated in quarter-length units, and
carries no information beyond it. One consequence reaches the reference project
directly: its \texttt{margin\_rank} state family --- eight features built from
``99.8\%-covered parsed \texttt{LBW} histories'' --- and its
\texttt{dynamic\_speed} family are measuring the same physical quantity, and any
apparent complementarity between them is quantisation noise.

\subsection{Q1: which target}

Targets are compared on \emph{forward persistence}: does a horse's date-atomic
prior mean predict its next realisation. A target a horse cannot repeat carries no
ability signal to find.

\begin{table}[htbp]\centering\small
\caption{Forward persistence of each candidate target, $n=104{,}539$.}
\begin{tabular}{lrrr}
\toprule
Target & Forward $r$ & Forward $R^2$ & Split-half reliability\\
\midrule
$y^{\text{rel}}_z$ \ within-race $z$ of log speed & \textbf{0.2777} & \textbf{0.0771} & 0.822\\
$y^{\text{rank}}$ \ within-race rank percentile & 0.2643 & 0.0698 & 0.840\\
$y^{\text{par}}$ \ par residual, log speed & 0.2376 & 0.0565 & 0.557\\
$y^{\text{par}}_z$ \ standardised par residual & 0.2357 & 0.0556 & 0.579\\
$y^{\text{par}}$ \ par residual, seconds & 0.2195 & 0.0482 & 0.460\\
$y^{\text{beaten}}$ \ beaten lengths & 0.2132 & 0.0454 & 0.431\\
\bottomrule
\end{tabular}
\end{table}

Within-race normalisation wins, and the reason is structural: it removes the
race-level shock, which is 32\% of the absolute residual's variance and is pure
noise for the purpose of estimating a horse's ability. Both targets are carried
forward, because the absolute one is the only one that can answer a question about
absolute environmental effects.

\figref{f04_targets}{Model 0 and the target. Left: incremental $R^2$ of log speed
by environment block. Centre: beaten seconds against beaten lengths --- the fit is
0.15996\,s per length, not the conventional 0.2, and the relationship is
arithmetic rather than statistical. Right: forward persistence of each candidate
target.}

\newpage
%======================================================================
\section{Exploratory analysis: what is identifiable}
%======================================================================

\begin{table}[htbp]\centering\small
\caption{History depth per estimation unit. This table determines what the rest of
the study can and cannot answer.}
\begin{tabular}{lrrrrrr}
\toprule
Unit & Units & Median & $<5$ & $<10$ & $\ge 20$ & Max\\
\midrule
Starts per horse & 8{,}599 & 15 & 16.3\% & 33.0\% & 41.0\% & 97\\
Rides per jockey & 190 & 121 & 19.5\% & 27.9\% & 62.6\% & 9{,}980\\
Rides per horse--jockey pair & 68{,}454 & \bad{1} & \bad{87.8\%} & 97.5\% & 0.23\% & 45\\
\bottomrule
\end{tabular}
\end{table}

The third row is decisive. A horse--jockey interaction estimated from a median of
one observation per cell is not an estimate, and \S\ref{sec:pairs} confirms it
formally rather than asserting it.

Horse $\times$ environment cells are thin but not hopeless: the median
horse--going cell holds 3 starts and 59.6\% hold fewer than five, while the median
horse--venue cell holds 8 and 42.2\% reach ten or more. That difference is why
\S\ref{sec:horseenv} finds a persistent venue sensitivity and no persistent going
sensitivity, and it is a difference in \emph{measurability} as much as in nature.

\figref{f01_sparsity}{History depth per estimation unit. The right-hand panel is
the constraint on Q4.}

\figref{f02_interaction_sparsity}{Why partial pooling is not optional. Left:
the fraction of interaction cells reaching each size. Right: one well-raced
horse's record across going levels.}

\newpage
%======================================================================
\section{Horse and jockey latent effects}
%======================================================================

\subsection{Estimation}

Crossed random effects are estimated by EM on the mixed-model equations: the E
step is the best linear unbiased predictor at the current variance ratios, solved
by Gauss--Seidel backfitting with \texttt{np.bincount}, and the M step is the
variance-component update that adds each effect's posterior variance to its
squared mean. Omitting that posterior-variance term is the classic error and
biases every $\tau^2$ downward.

The estimator is verified rather than trusted. On simulated crossed data with
$\tau_{\text{horse}}=0.55$, $\tau_{\text{jockey}}=0.20$, $\sigma=0.85$ it recovers
$0.544$, $0.192$, $0.851$. On data where the jockey factor has \emph{no} effect it
returns $\tau_{\text{jockey}}=0.0093$, a variance share of $8.5\times10^{-5}$ ---
so an effect it reports is not an artifact of the machinery.

\subsection{Q2 and Q3: variance components}

\begin{table}[htbp]\centering\small
\caption{Variance shares on within-race relative performance.}
\begin{tabular}{lrrrr}
\toprule
Model & Horse & Jockey & Trainer & Residual\\
\midrule
1: horse & 23.3\% & --- & --- & 76.7\%\\
2: $+$ jockey & 18.4\% & 3.7\% & --- & 78.0\%\\
2b: $+$ trainer & 17.7\% & 3.5\% & 1.2\% & 77.7\%\\
3: $+$ observable covariates & 17.6\% & 3.0\% & 1.1\% & 78.3\%\\
\midrule
Jockey alone, no horse control & --- & 6.0\% & --- & 94.0\%\\
Trainer alone & --- & --- & 1.4\% & 98.6\%\\
\bottomrule
\end{tabular}
\end{table}

Controlling for the horse removes about \textbf{39\%} of the apparent jockey
variance (6.0\% to 3.7\%). That is the size of the assignment channel: good
jockeys ride good horses.

\subsection{The surprise: jockey identity forecasts better than horse identity}

\begin{table}[htbp]\centering\small
\caption{Forward evaluation on the confirmation window, 2020--2026. The scale is
fitted on the previous year's own earlier fit, so it never sees the scored year.}
\begin{tabular}{lrrrr}
\toprule
Model & oos $R^2$ (scaled) & Forward scale & Mean $r$ & Years positive\\
\midrule
1: horse & 0.0124 & 0.39 & 0.141 & 7/7\\
\textbf{Jockey alone} & \good{0.0479} & \good{1.05} & 0.223 & 7/7\\
2: horse $+$ jockey & 0.0496 & 0.66 & 0.227 & 7/7\\
2b: $+$ trainer & 0.0536 & 0.67 & 0.236 & 7/7\\
3: $+$ covariates & \textbf{0.0582} & 0.64 & 0.245 & 7/7\\
Trainer alone & 0.0063 & 0.85 & 0.094 & 7/7\\
\bottomrule
\end{tabular}
\end{table}

Jockey identity alone forecasts nearly four times as well as horse identity, and
at a forward scale of 1.05 --- essentially perfectly calibrated --- where the horse
effect needs shrinking to 0.39 of its fitted size. Adding the horse to the jockey
improves oos $R^2$ from 0.0479 only to 0.0496.

The explanation is precision and drift, not size. A jockey effect rests on roughly
a thousand rides and does not move; a horse effect rests on fifteen starts and
does. Two consequences worth separating:

\begin{itemize}\itemsep2pt
\item As a \emph{predictor}, jockey identity is powerful partly \emph{because} it
encodes assignment --- the market knows this too, which is why the production
model's \texttt{jockey\_residual} is a residual \emph{against the market} rather
than a strike rate.
\item As a \emph{causal} effect it is much smaller. The horse-adjusted share is
3.7\%, and the plan's own instruction not to read a jockey random effect causally
without discussing assignment bias is exactly right.
\end{itemize}

\subsection{Identification diagnostics}

Because jockey assignment is highly non-random, the effect is estimated a second,
independent way: from \emph{within-horse} variation only, using each ride's
deviation from that horse's own mean.

The within-horse estimates have standard deviation 0.141 against a median standard
error of 0.028 --- \good{5.05 times their own noise}, so they are not sampling
error --- and they agree with the crossed BLUPs at Spearman \good{0.970}. Horse
control shrinks the spread of jockey effects by 17.7\%, and nine of the top ten
jockeys by raw average remain in the top ten after adjustment.

\subsection{Q2: ability drifts, with a two-month half-life}

\begin{table}[htbp]\centering\small
\caption{Forward correlation of a decayed horse-ability estimate. Correlation is
the criterion because it is scale-free; the optimal scale column shows how badly
a static estimate is over-scaled.}
\begin{tabular}{lrrr}
\toprule
Recency half-life & Forward $r$ & $R^2$ at best scale & Implied optimal scale\\
\midrule
30 days & 0.1445 & 0.00065 & 0.870\\
45 days & 0.1485 & 0.00244 & 0.772\\
\textbf{60 days} & \good{0.1497} & \textbf{0.00314} & 0.716\\
90 days & 0.1488 & 0.00325 & 0.650\\
180 days & 0.1407 & 0.00106 & 0.550\\
365 days & 0.1301 & $-0.00201$ & 0.470\\
730 days & 0.1227 & $-0.00406$ & 0.422\\
Static, all history & \bad{0.1142} & $-0.00628$ & \bad{0.369}\\
\bottomrule
\end{tabular}
\end{table}

The ordering is monotone with no interior structure to over-read: recent form
matters far more than distant form, and the whole career is the worst estimator
tested. This speaks directly to the reference project's open backlog item on
recency-weighted training, where truncating history was found to hurt badly
overall while helping recent years --- weighting rather than truncating is the
untried alternative, and the optimum here is a 60-day half-life.

\figref{f08_variance_components}{Left: variance shares by model. Centre: the
jockey effect before and after horse control, point size proportional to rides.
Right: forward correlation against recency half-life; the static estimate is the
worst of the eight tested.}

\figref{f09_horse_ability}{Shrunk horse ability against starts, and the posterior
standard deviation of the estimate. The right-hand panel is a quantity both of the
reference project's latent engines compute and neither promotes.}

\newpage
%======================================================================
\section{Environmental effects: absolute time against relative ordering}
%======================================================================

The plan's Q5 asks for both quantities separately, because they are routinely
confused and they point in opposite directions here.

\textbf{On absolute time, the environment dominates.} It explains 0.9962 of the
variance of raw seconds. Holding distance, venue and surface fixed at Sha Tin turf
1200\,m, going alone moves mean finishing time by \textbf{1.196\,s} across
adequately-supported levels --- roughly seven and a half lengths.

\textbf{On relative ordering, almost nothing survives.} After the forward par
model, the residual race-level shock still has a standard deviation of 0.00724 in
log speed (about half a second at 1200\,m), and the observable card explains
almost none of it:

\begin{table}[htbp]\centering\small
\caption{One-way ANOVA $R^2$ of the residual race-level shock, count-weighted.}
\begin{tabular}{lrrr}
\toprule
Dimension & Levels ($\ge 200$ races) & ANOVA $R^2$ & Range in seconds at 1200\,m\\
\midrule
going & 3 & 0.0298 & 0.221\\
venue & 2 & 0.0099 & 0.102\\
distance band & 4 & 0.0058 & 0.116\\
class & 4 & 0.0008 & 0.050\\
course letter & 4 & 0.0005 & 0.035\\
surface & 2 & 0.0004 & 0.032\\
\bottomrule
\end{tabular}
\end{table}

So roughly \textbf{96\% of the race-level shock is unobserved} from the card ---
day-to-day track speed, weather and tactics. \S\ref{sec:race} shows that most of
it \emph{is} observable, just not from the card: it is observable from the rest of
the meeting.

\figref{f14_environment_absolute}{Left: going against absolute finishing time on
one distance, track and surface. Right: draw against the win rate and against the
market's price for the same runners.}

\newpage
%======================================================================
\section{Horse \texorpdfstring{$\times$}{x} environment effects}\label{sec:horseenv}
%======================================================================

The plan calls this the single most important question, and prescribes the right
test: what matters is not a horse's apparent sensitivity but the between-horse
variance of the \emph{true} sensitivity, and the only evidence for that is whether
an early estimate predicts a later realisation.

The design is deliberately hard to fool. For each binary environmental contrast,
form each horse's performance contrast from its \emph{early} starts; decompose the
spread of those estimates into true between-horse variance and sampling noise by
method of moments; then test the early estimate against the same horse's
\emph{later} contrast.

\begin{table}[htbp]\centering\small
\caption{Horse $\times$ environment sensitivity. ``Real share'' is the fraction of
the apparent spread that is not sampling noise. Persistence is the correlation
between the early and later estimates, with $z$ against zero.}
\begin{tabular}{lrrrrr}
\toprule
Contrast & Horses & Apparent sd & Real $\tau$ & Real share & Persistence\\
\midrule
going, soft vs firm & 1{,}675 & 0.515 & 0.216 & 0.175 & \bad{$+0.034$ ($z=1.4$)}\\
\textbf{venue, HV vs ST} & 1{,}476 & 0.599 & 0.374 & 0.389 & \good{$+0.215$ ($z=8.3$)}\\
\textbf{surface, AWT vs turf} & 365 & 0.607 & 0.402 & 0.438 & \good{$+0.338$ ($z=6.4$)}\\
distance, long vs short & 163 & 0.634 & 0.387 & 0.374 & \bad{$+0.064$ ($z=0.8$)}\\
field size, large vs small & 36 & 0.608 & 0.306 & 0.254 & \bad{$-0.138$ ($z=-0.8$)}\\
class, high vs low & --- & --- & --- & --- & unanswerable\\
\bottomrule
\end{tabular}
\end{table}

Three readings, in order of importance.

\textbf{Most apparent sensitivity is noise, everywhere.} Even for the contrasts
that persist, the real between-horse standard deviation is 0.37--0.40 against an
apparent 0.60 --- so 56--62\% of what a naive per-horse table would show is
sampling error. A study that stopped at the apparent spread would have reported
five real effects.

\textbf{Going sensitivity does not persist.} This is the plan's own nominated
``most valuable possible finding'' and it is confirmed: going moves absolute race
time by 1.2\,s, and horse-specific going sensitivity has a forward correlation of
0.034 at $z=1.4$ on 1{,}675 horses --- the largest sample of the five contrasts.
The effect on \emph{everyone} is enormous; the differential effect is not
measurable, and is not usefully there.

\textbf{Venue and surface sensitivity are real.} ``Horses for courses'' between
Happy Valley's tight circuit and Sha Tin's galloping track survives at $z=8.3$
with a true between-horse standard deviation of 0.374 performance standard
deviations --- a materially large effect. Surface sensitivity is stronger still
($z=6.4$, $\tau=0.402$) on a much smaller qualifying sample.

The measurability difference tracks the cell counts exactly: the median
horse--venue cell has 8 starts against the median horse--going cell's 3. Part of
this result is about nature and part is about what 13{,}309 races can resolve, and
the honest statement separates them.

\figref{f13_environment_sensitivity}{Left: apparent against real between-horse
spread for each contrast --- the gap is sampling noise. Centre: persistence with
its null standard error. Right: early against later sensitivity for the strongest
contrast, binned.}

\newpage
%======================================================================
\section{The residual distribution}
%======================================================================

Measured on 96{,}579 genuinely out-of-sample residuals from a forward-fitted mean
model with the forward scale applied, 2016--2026.

\begin{table}[htbp]\centering\small
\caption{Shape of the standardised out-of-sample residual.}
\begin{tabular}{lrr}
\toprule
Statistic & Value & Normal\\
\midrule
Skewness & \bad{$-0.729$} & 0\\
Excess kurtosis & $+0.810$ & 0\\
IQR / sd & 1.1195 & 1.3490\\
Mean $-$ median & $-0.070$ & 0\\
\bottomrule
\end{tabular}
\end{table}

The centre is \emph{more} concentrated than a normal and the tails are heavier ---
IQR/sd of 1.12 against 1.35 --- with a pronounced left skew. Candidate densities
were fitted on the development window and scored on the confirmation window by
predictive log-likelihood, so no family is preferred for fitting its own data.

\begin{table}[htbp]\centering\small
\caption{Held-out predictive log-likelihood per observation, relative to the Normal.}
\begin{tabular}{lrr}
\toprule
Family & Parameters & Gain over Normal per observation\\
\midrule
Two-component normal mixture & 5 & \good{$+0.0557$}\\
Skew-$t$ & 4 & $+0.0553$\\
Skew-normal & 3 & $+0.0455$\\
Student-$t$ ($\nu=5.18$) & 3 & $+0.0246$\\
Normal & 2 & 0\\
\bottomrule
\end{tabular}
\end{table}

\subsection{Which tail, and what it was}

The plan asks whether the bad-performance tail is materially heavier, and whether
the outliers are identifiable events. Both answers are yes, and the second is
checkable because the reference archive holds the stewards' text.

\begin{table}[htbp]\centering\small
\caption{Stewards'-comment rates by residual band. ``Severe'' matches the language
used for a run that was stopped or abandoned; its base rate is 9.5\%.}
\begin{tabular}{lrrrr}
\toprule
Residual band & $n$ & Any-trouble rate & Severe rate & Severe lift\\
\midrule
$<-3$ sd & 889 & 0.693 & \good{0.460} & \good{4.83}\\
$-3$ to $-2$ & 4{,}200 & 0.624 & 0.259 & 2.71\\
$-2$ to $-1$ & 10{,}367 & 0.608 & 0.157 & 1.65\\
$-1$ to $+1$ & 70{,}122 & 0.575 & 0.078 & 0.82\\
$+1$ to $+2$ & 10{,}240 & 0.501 & 0.057 & 0.59\\
$+2$ to $+3$ & 755 & 0.437 & 0.050 & 0.53\\
\bottomrule
\end{tabular}
\end{table}

The gradient is monotone across all bands in both tiers. Nearly half of the runs
three standard deviations below expectation carry a severe stewards' phrase, at
4.8 times the base rate --- and good performances are \emph{less} likely than
average to carry one.

Two interpretive points the plan explicitly asks for. First, these observations
are \emph{not} removed: a disastrous run is part of the outcome distribution a
probability model must represent, and the reference project's own settlement
treats non-finishers as losing starters for the same reason. Second, part of the
left skew is \textbf{structural rather than stochastic}: a beaten horse is eased,
so the lower tail of a within-race time distribution is long by the physics of the
sport. The mixture's fitted ``bad'' component carries probability 0.44, which is
far too large to read as a bad-trip rate and is better read as the asymmetry of
the racing outcome itself.

\figref{f10_residual_distribution}{Left: the residual on a log density scale with
Normal and Student-$t$ overlays. Centre: the Q--Q plot. Right: empirical over
Normal tail probabilities, split by sign --- the bad tail is the heavy one.}

\figref{f11_residual_families}{Left: held-out log-likelihood gain over the Normal.
Right: stewards'-comment rates by residual band, two severity tiers.}

\newpage
%======================================================================
\section{Heteroskedasticity and horse-specific volatility}
%======================================================================

The plan judges constant variance more important than exact normality, and it is
right: a model with the mean right and the spread wrong is miscalibrated where
race log loss punishes hardest.

\subsection{Q9: the variance is not constant}

\begin{table}[htbp]\centering\small
\caption{Residual spread by decile of each covariate.}
\begin{tabular}{lrrr}
\toprule
Covariate & sd, lowest decile & sd, highest decile & Ratio\\
\midrule
Market rank & 0.812 & 1.051 & \textbf{1.295}\\
Market probability & 1.051 & 0.824 & 0.784\\
Draw position & 0.914 & 0.984 & 1.077\\
Field size & 0.933 & 0.949 & 1.017\\
Carried weight & 0.960 & 0.936 & 0.975\\
\bottomrule
\end{tabular}
\end{table}

The largest driver is the market's own assessment, and the direction is
informative: \textbf{favourites run more consistently}. A joint model of
$\log|e|$ on observables reaches out-of-sample $R^2$ 0.011 with an implied
p90/p10 spread ratio of 1.385.

\subsection{Q10: volatility is a horse trait, and it persists}

The naive per-horse standard deviation is dominated by sample size, exactly as the
plan requires be demonstrated. Decomposing the observed variance of per-horse
$\log$ sd:

\begin{table}[htbp]\centering\small
\caption{Variance components of per-horse log residual standard deviation,
$n=5{,}324$ horses with at least two residuals.}
\begin{tabular}{lr}
\toprule
Quantity & Value\\
\midrule
Observed variance of $\log \hat{\sigma}_i$ & 0.2060\\
Mean sampling variance $\approx 1/(2(n_i-1))$ & 0.0786\\
True between-horse variance $\tau^2$ & 0.1274\\
\textbf{Share of the observed spread that is real} & \good{61.8\%}\\
$\tau$ (sd of true $\log\sigma$) & 0.357\\
Spread ratio for a horse 1 sd more volatile & \good{1.43$\times$}\\
\bottomrule
\end{tabular}
\end{table}

The decisive test is the plan's own: split each horse's residual history at its
own median date and ask whether the first half's spread predicts the second's.

\begin{table}[htbp]\centering\small
\caption{Split-half persistence of per-horse log residual sd.}
\begin{tabular}{lrrrr}
\toprule
Min starts each half & Horses & $r$ & $z$ vs zero & Later-sd ratio, top/bottom tercile\\
\midrule
4 & 3{,}876 & $+0.144$ & $+8.94$ & 1.115\\
6 & 3{,}146 & $+0.156$ & $+8.77$ & 1.119\\
8 & 2{,}519 & $+0.176$ & $+8.85$ & 1.131\\
10 & 2{,}015 & $+0.170$ & $+7.63$ & 1.115\\
15 & 1{,}024 & \good{$+0.223$} & $+7.14$ & 1.148\\
\bottomrule
\end{tabular}
\end{table}

The correlation rises with depth, which is what a real trait measured with less
noise should do. As a forward feature, a prior-only volatility estimate predicts
the next start's $|residual|$ at Spearman $+0.087$ to $+0.122$, with the realised
spread ratio between the top and bottom terciles reaching 1.255 at twelve prior
starts.

\textbf{So horse-specific volatility passes every test the plan sets for it.}
Whether it is \emph{useful} is a separate question, and \S\ref{sec:generative} and
\S\ref{sec:features} answer it differently --- which is the most instructive
disagreement in this study.

\figref{f12_volatility}{Left: naive against shrunk per-horse volatility by
observation count, with the pooled value. Centre: split-half persistence with its
null standard error. Right: the covariates that move the spread.}

\newpage
%======================================================================
\section{Horse \texorpdfstring{$\times$}{x} jockey compatibility}\label{sec:pairs}
%======================================================================

\subsection{A correction, recorded}

The first version of this analysis reported early-to-late pair correlations of
$-0.04$ to $-0.21$ with $z$ from $-3.8$ to $-7.4$ --- an apparently strong
\emph{negative} result that would have read as compatibility mean-reverting.

It was an artifact of the residualisation, and the mechanism is worth stating
because it is easy to reproduce. The horse and jockey main effects were fitted
over the whole sample, so a horse's effect is roughly the average of its early and
late performance. The early pair residual then becomes about
$(\text{early}-\text{late})/2$ and the late one about
$(\text{late}-\text{early})/2$: forced anticorrelation with no compatibility
involved. The artifact was \emph{larger in magnitude} than a true pair standard
deviation of 0.30 would have produced, which is how it was caught.

The corrected design fits the main effects separately within each era, so the
early estimate and the late realisation share no fitted parameter.

\subsection{The corrected answer, with power}

\begin{table}[htbp]\centering\small
\caption{Forecastability of a pair effect, with the correlation a real effect of
each size would have produced on the same design.}
\begin{tabular}{lrrrrr}
\toprule
Min rides each half & Pairs & Apparent sd & True $\tau$ & $r$ & $z$\\
\midrule
2 & 575 & 0.483 & 0.234 & $-0.002$ & $-0.05$\\
3 & 249 & 0.383 & 0.155 & $+0.024$ & $+0.38$\\
4 & 121 & 0.303 & \textbf{0.000} & $+0.071$ & $+0.77$\\
5 & 61 & 0.219 & \textbf{0.000} & $-0.053$ & $-0.41$\\
\midrule
\multicolumn{6}{l}{\emph{Power: what a real pair effect would have shown}}\\
\multicolumn{4}{l}{true pair sd $=0.10$} & $+0.058$ & \\
\multicolumn{4}{l}{true pair sd $=0.20$} & $+0.211$ & \\
\multicolumn{4}{l}{true pair sd $=0.30$} & $+0.334$ & \\
\bottomrule
\end{tabular}
\end{table}

Meanwhile the \textbf{in-sample} Model 7 fit, on pairs with three or more rides,
reports $\tau_{HJ}=0.298$ and a \bad{10.3\%} variance share.

That contrast is the finding. An in-sample pair effect apparently worth a tenth of
the variance, and a forecastable component indistinguishable from zero on a design
with the power to see half that size. Q4 is answered negatively, and the negative
is supported rather than merely reported.

\figref{f16_pairs}{Left: rides per horse--jockey pair on a log count scale ---
52\% ride once. Right: measured persistence against what a real effect of each
size would have produced.}

\newpage
%======================================================================
\section{Shared race effects and pace}\label{sec:race}
%======================================================================

The reference project already tried a one-factor structured-covariance probit and
got a negative result, with the factor scale essentially unidentified: its grid
selected ``no factor'' in seven of eleven years and the whole range from 0.00 to
0.80 spanned about $2\times10^{-5}$ of log loss. Two features of that test left the
question open. Its loading was a running-\emph{position} habit parsed from the
result page's three-number line rather than a timed pace measurement, and its
standardisation used the whole frame including the test block.

\subsection{The shock is large, and mostly a property of the day}

The shared shock is \textbf{32.1\%} of the absolute residual's variance, with a
standard deviation of 0.00724 in log speed --- about \textbf{0.503\,s at 1200\,m}.

\begin{table}[htbp]\centering\small
\caption{How predictable a race's own shock is.}
\begin{tabular}{lr}
\toprule
Predictor & Correlation with the race's shock\\
\midrule
Running mean of earlier races on the same card & \good{0.504}\\
\quad by the 8th race of the meeting & \good{0.584}\\
Prior history of the same venue and surface & 0.179\\
\bottomrule
\end{tabular}
\end{table}

Track speed is therefore a day-level property and it is largely \emph{measurable}
before a later race on the card runs. This is the one signal in this study that is
both large and structurally unavailable to the reference pipeline, whose
point-in-time rule excludes all same-date information. That rule exists to stop a
horse seeing its own same-day outcome; a track-speed reading taken from race one
before race eight goes to post is forward in time and is not that. It is recorded
here as a labelled deviation and is not folded into any candidate feature.

\subsection{Q11: the shock does load differently by running style}

A shock loading equally on every runner cancels from every pairwise utility
difference and cannot change any probability, so the interaction is the whole
question. Measuring the race's pace directly from the sectional archive:

\begin{table}[htbp]\centering\small
\caption{Pace-by-style interactions on within-race relative performance,
development window fitted, confirmation window scored.}
\begin{tabular}{lrrr}
\toprule
Specification & Interaction coefficient & $z$ & oos $R^2$ gain\\
\midrule
Forwardness $\times$ early pace & $-0.00034$ & \good{$-8.31$} & $-0.00002$\\
Forwardness $\times$ pace shape & $+0.00036$ & $+8.55$ & $+0.00035$\\
Closing ability $\times$ early pace & $+0.00031$ & $+7.43$ & $+0.00118$\\
\bottomrule
\end{tabular}
\end{table}

The signs are the physics: a fast early pace \emph{hurts} habitual front-runners
and \emph{helps} closers. So the mechanism the reference project's one-factor
probit hypothesised is real and correctly signed --- it was undetectable with a
loading that could not see pace. Its economic size is small.

\subsection{Is the pace forecastable, and does the mispricing survive?}

Everything above conditions on the pace the race actually ran. A model can only
use a pace it can project.

\begin{table}[htbp]\centering\small
\caption{Out-of-sample forecastability of a race's early pace.}
\begin{tabular}{lrr}
\toprule
Projection & oos $r$ & oos $R^2$\\
\midrule
Top-3 prior style only & 0.032 & $-0.005$\\
\textbf{Field style composition (12 features)} & \good{0.414} & \good{0.158}\\
$+$ earlier races on the same card & 0.405 & 0.153\\
Earlier races on the same card only & 0.098 & $-0.003$\\
\bottomrule
\end{tabular}
\end{table}

A naive summary sees almost nothing; a proper projection recovers $r=0.41$. Note
also that track speed and early pace are largely \emph{distinct} quantities ---
card track speed alone predicts early pace at only 0.098 --- so the surface's
speed and how hard they went early are separate facts.

The market's error across the style $\times$ pace grid swings $-0.0398$ on the
\emph{realised} pace, and $-0.0151$ on the \emph{projected} pace: \good{37.8\%
survives}, with a maximum absolute residual of 0.0087 on 50{,}784 confirmation-window
runners. Roughly nine tenths of a percentage point of win probability, available
pre-race, at the extremes of the grid.

\figref{f15_race_factors}{Left: the distribution of the race-level shock. Centre:
relative performance by prior style and realised early pace --- the interaction is
visible. Right: the market residual over the same grid.}

\newpage
%======================================================================
\section{From latent performance to probabilities}\label{sec:generative}
%======================================================================

\subsection{Integration, and why not simulation}

For independent shocks the winner probability collapses to a one-dimensional
integral over the winner's own shock,
\[
P(i\ \text{wins}) = \mathbb{E}_z\!\left[\prod_{j\ne i} F_j\!\left(\frac{\mu_i-\mu_j+\sigma_i z}{\sigma_j}\right)\right],
\]
which holds for \emph{heterogeneous} $\sigma$ as written. Simulation is avoided for
a reason the reference project measured: a winner-counting estimator over $m$ draws
cannot resolve a probability below $1/m$, and this sport's longshots sit at
$10^{-3}$ to $10^{-4}$ where race log loss is decided.

The integrator is verified against every case with a known answer: two runners at
equal scale to $1.1\times10^{-16}$, at unequal scale to $1.3\times10^{-12}$, an
identical twelve-runner field to $1.4\times10^{-17}$ from uniform, a longshot
resolved at $2.3\times10^{-5}$, total mass to $1.1\times10^{-16}$, a constant
shared-factor loading collapsing onto the iid answer exactly, and the Student-$t$
path against a 400{,}000-draw simulation to within its simulation standard error.

\subsection{Q13: what each piece of structure buys}

Six nested variants, each a controlled comparison against the row above. One
global utility scale is fitted per variant on the last development year and applied
forward; that year is scored on nothing.

\begin{table}[htbp]\centering\small
\caption{Nested generative variants. 59{,}364 runners in 4{,}844 races, 2020--2026.
Market race log loss 2.02718.}
\begin{tabular}{lrrrr}
\toprule
Variant & Fitted scale & Race log loss & vs market & Gain over previous\\
\midrule
A. homoskedastic normal & 0.73 & \good{2.39671} & $+0.36953$ & ---\\
B. $+$ heteroskedastic (observables) & 0.84 & 2.40665 & $+0.37947$ & \bad{$-0.00994$}\\
C. $+$ horse-specific volatility & 0.98 & 2.41790 & $+0.39072$ & \bad{$-0.01125$}\\
D. $+$ Student-$t$ shocks & 1.24 & 2.40726 & $+0.38008$ & $+0.01064$\\
E. $+$ latent-state uncertainty & 1.30 & 2.39776 & $+0.37058$ & $+0.00950$\\
F. $+$ shared race factor on style & 1.30 & 2.39947 & $+0.37229$ & \bad{$-0.00171$}\\
\bottomrule
\end{tabular}
\end{table}

\textbf{No elaboration beats plain homoskedastic normal shocks.} This is the
central negative result of the study and it survived a correction: an earlier run
of this stage left the utility scale at 1.0 through a bug and reported that
heteroskedasticity, $t$ tails and latent-state uncertainty all helped. Fitting the
scale per variant reversed the conclusion. The mechanism is visible in the scale
column: each elaboration changes the relative spread of runners' shocks, and the
fitted global scale then absorbs most of what was gained.

This independently reproduces the reference project's own probit finding that the
utility scale is ``the whole game'' --- now with measured heteroskedasticity and a
sectional-derived pace loading rather than a position habit, so it is a stronger
statement of the same result.

One further observation. Variant A's calibration RMSE is \good{0.00402} against
the market's 0.00524 and variant F's 0.00635: the simplest generative model is
\emph{better calibrated than the market} while being 0.37 of log loss worse. It is
well calibrated and far less sharp, which is exactly what a market-blind latent
model should look like, and it is a reminder that calibration alone is not a
measure of value.

\subsection{Q12: reducible against irreducible uncertainty}

\begin{table}[htbp]\centering\small
\caption{The predictive variance budget by the horse's prior starts. ``Reducible''
is the share attributable to not yet knowing the horse's ability.}
\begin{tabular}{lrrrr}
\toprule
Prior starts & $n$ & Latent-state sd & Idiosyncratic sd & Share reducible\\
\midrule
0--3 & 13{,}203 & 0.409 & 0.845 & \good{19.1\%}\\
4--8 & 13{,}527 & 0.367 & 0.846 & 15.8\%\\
9--15 & 14{,}502 & 0.259 & 0.837 & 8.4\%\\
16--30 & 18{,}129 & 0.195 & 0.850 & 5.0\%\\
30$+$ & 9{,}263 & 0.147 & 0.865 & \good{2.8\%}\\
\bottomrule
\end{tabular}
\end{table}

Clean and monotone. The idiosyncratic component is essentially constant at
0.84--0.87 across all five bands, while the reducible part falls from 19\% at
three starts to 3\% at thirty. \textbf{Beyond about fifteen starts, nearly all
remaining uncertainty is irreducible conditional performance randomness}, and no
amount of extra history will help.

The shared race component enters the predictive variance of \emph{absolute}
performance and barely touches the win probability, because it is common to the
field and cancels from the ordering except through the style loading.

\figref{f17_probabilities}{Left: what each nested structure buys. Centre:
calibration on log--log axes for the simplest variant, the full one and the
market. Right: the reducible share of predictive variance against history depth.}

\figref{f18_model_vs_market}{The generative model's probability against the
market's, on log axes. The dispersion is what a 0.37 log-loss gap looks like.}

\newpage
%======================================================================
\section{Market comparison and feature integration}\label{sec:features}
%======================================================================

This is the plan's declared primary practical output, and the only stage whose
metric is the one a promotion decision turns on.

\subsection{Method, and why the numbers are trustworthy}

The production offset architecture is reimplemented independently
(\texttt{src/hkjcres/offset.py}) rather than imported, so this experiment cannot
mutate the reference project and the architecture has to be stated rather than
assumed. It verifies: $\operatorname{softmax}(\log q)$ reproduces $q$ to
$1.4\times10^{-16}$, and a zero-round fit reproduces the market to
$4.1\times10^{-8}$.

Three disciplines are taken from the reference project:

\begin{itemize}\itemsep2pt
\item The arms share everything by construction --- one split per year, one label
vector, one base margin, one seed, computed once and handed to both fits, with the
candidate's column list being the incumbent's with the block \emph{appended},
because XGBoost's column indices are positional.
\item \textbf{The no-op candidate must return exactly zero.} It does: the maximum
absolute cell delta across fifteen cells is \good{0.0}.
\item No dividends are loaded and no return is computed. At this effect size ROI
moves by tens of percentage points on changes that leave log loss unchanged.
\end{itemize}

Development profile: test years 2016, 2019, 2021, 2024, 2025; seeds 7, 17, 42;
150 rounds; 122{,}880 rows over 9{,}934 races. On this profile the incumbent beats
the market by $-0.00109$.

\subsection{The ranked candidate table}

\begin{table}[htbp]\centering\small
\caption{Candidate blocks against the production incumbent. Negative delta is
better. \emph{Novel share} is one minus the out-of-sample $R^2$ with which the
incumbent's own 28 columns reconstruct the candidate.}
\begin{tabular}{lrrrrrr}
\toprule
Block & Mean $\Delta$ & s.e. & $t$ & Cells & Years & Novel share\\
\midrule
Full sectional profile (Q1B) & $-0.001179$ & 0.00130 & $-0.91$ & 9/15 & 3/5 & 0.51--0.83\\
Recent-form ability (Q2) & $-0.001018$ & 0.00106 & $-0.96$ & 8/15 & 3/5 & 0.39\\
Sectional finish (Q1B) & $-0.000870$ & 0.00081 & $-1.07$ & 8/15 & 3/5 & 0.52\\
Pace fit (Q11) & $-0.000283$ & 0.00141 & $-0.20$ & 8/15 & 2/5 & 0.76\\
Horse volatility (Q10) & $-0.000110$ & 0.00087 & $-0.13$ & 9/15 & 2/5 & 0.90\\
\rowcolor{shade} \emph{no-op control} & \emph{0.000000} & --- & --- & --- & --- & ---\\
Ability uncertainty (Q12) & $+0.000050$ & 0.00059 & $+0.09$ & 6/15 & 2/5 & 0.27\\
Sectional middle (Q1B, new) & $+0.000686$ & 0.00071 & $+0.97$ & 6/15 & 2/5 & 0.83\\
Everything that survived & $+0.001167$ & 0.00155 & $+0.75$ & 6/15 & 1/5 & ---\\
Condition fit (Q6) & $+0.001379$ & 0.00105 & $+1.31$ & 4/15 & 2/5 & 0.93\\
\bottomrule
\end{tabular}
\end{table}

\textbf{No block clears a promotion gate.} The best three are directionally
favourable at $-0.00087$ to $-0.00118$ --- 20 to 26\% of the production model's
entire edge over the market --- but with $|t|\approx 1$, three of five years
improving, and a standard error of the same order as the mean. On the reference
project's own criteria (sign per year, worst year, sign per seed, then the mean)
that is \emph{investigate}, not \emph{promote}.

\subsection{Q14: what the market has not priced}

Separately from model value, the plan asks whether each quantity predicts
\emph{market error}. A linear probability model of the outcome on the market's own
logit plus the standardised candidate:

\begin{table}[htbp]\centering\small
\caption{Does each candidate explain the outcome given the market's logit?}
\begin{tabular}{lrrr}
\toprule
Feature & Coefficient & $t$ & Correlation with the market price\\
\midrule
Ability uncertainty & $+0.0113$ & \good{$+14.99$} & $-0.004$\\
History depth & $-0.0110$ & $-14.65$ & $-0.000$\\
Sectional middle & $-0.0047$ & $-6.14$ & $+0.169$\\
Projected pace & $+0.0041$ & $+5.51$ & $-0.035$\\
Sectional finish & $-0.0031$ & $-3.91$ & $+0.291$\\
Venue fit & $-0.0024$ & $-3.24$ & $+0.079$\\
Surface fit & $-0.0018$ & $-2.42$ & $+0.057$\\
Horse volatility & $-0.0017$ & $-2.23$ & $-0.102$\\
\bottomrule
\end{tabular}
\end{table}

\textbf{The market underprices lightly-raced horses.} A measure of ability
uncertainty essentially uncorrelated with the price ($-0.004$) predicts the
outcome given the price at $t=+15.0$, worth about 1.1 percentage points of win
probability per standard deviation. This is exactly the plan's Q14 hypothesis
about ``uncertainty for lightly raced horses'', and it is the strongest
market-inefficiency signal in the study.

\subsection{The resolution, and the study's most useful pattern}

That leaves a tension: a quantity with $t=15$ against the market adds
$+0.00005$ to the model. Regressing each candidate on the incumbent's own 28
columns, forward, resolves it.

\begin{table}[htbp]\centering\small
\caption{How much of each candidate the incumbent's 28 columns already imply,
out of sample.}
\begin{tabular}{lrr}
\toprule
Candidate & Boosted $R^2$ from the incumbent & Novel share\\
\midrule
Venue fit & 0.066 & 0.934\\
Horse volatility & 0.104 & 0.896\\
Surface fit & 0.111 & 0.890\\
Sectional middle & 0.166 & 0.834\\
Projected pace & 0.236 & 0.764\\
Sectional finish & 0.480 & 0.520\\
Recent-form ability & 0.608 & 0.392\\
History depth & 0.717 & 0.283\\
\textbf{Ability uncertainty} & \good{0.732} & 0.268\\
\bottomrule
\end{tabular}
\end{table}

The incumbent already reconstructs \textbf{73\%} of ability uncertainty from
\texttt{layoff\_log}, \texttt{bt\_days\_since}, \texttt{rating\_z},
\texttt{class\_change} and the rest. The market's inefficiency is real; the
production model has already captured it.

And across the nine blocks, the rank correlation between \emph{novel share} and
\emph{delta} is \good{$+0.389$} --- positive meaning \textbf{the more novel a
candidate, the less it helped}. The candidates that helped (sectional finish 0.52
novel, recent-form ability 0.39) are ones the model already partly had; the
genuinely novel ones (condition fit 0.93, horse volatility 0.90, sectional middle
0.83) did nothing or hurt, despite \S\ref{sec:horseenv} and \S9 establishing that
they are real and persistent.

The inference is uncomfortable and well-evidenced: at this effect size, marginal
value lies in \emph{sharpening information the model already uses}, not in adding
information it lacks. Orthogonal signal that is real but small is dominated by the
noise it introduces and by the tree capacity it consumes. This is a coherent
explanation for the reference project's own history --- roughly four hundred
candidate columns built, audited by expanding-year audits, and none promoted.

\figref{f19_features}{Left: candidate block deltas with standard errors. Centre:
market-residual $t$ statistics. Right: novelty against realised model value ---
the upward slope is the study's most useful single pattern.}

\newpage
%======================================================================
\section{Answers to the fourteen questions}
%======================================================================

\begin{longtable}{@{}p{0.055\textwidth}p{0.40\textwidth}p{0.50\textwidth}@{}}
\toprule
\textbf{Q} & \textbf{Question} & \textbf{Answer}\\
\midrule
\endhead
Q1 & Best observable proxy for latent performance & Within-race $z$ of log speed. Forward $R^2$ 0.077 against 0.056 for a par-time residual; reliability 0.822. Beaten lengths are arithmetically the time and carry nothing extra.\\[3pt]
Q1B & Do sectionals beat final time & \textbf{Yes for persistence, no for information.} Fourteen sectional quantities out-persist the best time target (best $r=0.483$ against 0.278), but the profile is \emph{not} low-dimensional (49.6/32.3/18.1\%) and the most persistent quantities are the least informative. Shape adds $+0.0026$ of oos $R^2$ over time alone.\\[3pt]
Q2 & How much variation is horse ability & 23.3\% of relative performance alone, 17.6\% after jockey, trainer and covariates. Ability \emph{drifts}: a 60-day half-life beats the full career (forward $r$ 0.150 against 0.114), and a static estimate is over-scaled 2.7$\times$.\\[3pt]
Q3 & Is jockey ability identifiable and useful & \textbf{Yes to both, and it forecasts better than horse ability.} 3.7\% of variance after horse control (from 6.0\% before, so 39\% of the apparent effect is assignment). Within-horse switches give a spread 5.05$\times$ their own noise and agree with the BLUPs at 0.970. Jockey alone: oos $R^2$ 0.048 at scale 1.05, against horse alone 0.012 at 0.39.\\[3pt]
Q4 & Horse $\times$ jockey interaction & \textbf{No.} Median pair has one ride. Corrected early-to-late correlation is null ($z$ from $-0.4$ to $+0.8$) on a design that would show $r=0.21$ for a true sd of 0.20. The in-sample fit's 10.3\% variance share is overfitting.\\[3pt]
Q5 & Environmental effects on absolute performance & Enormous. Environment explains 0.9962 of raw-seconds variance; going alone moves time 1.196\,s at ST turf 1200\,m. On \emph{relative} ordering, the observable card explains $\sim$3\% of the residual race shock.\\[3pt]
Q6 & Horse $\times$ environment heterogeneity & \textbf{Real for venue ($z=8.3$) and surface ($z=6.4$); not for going ($z=1.4$), distance or field size.} 56--62\% of even the real contrasts' apparent spread is sampling noise. Adding it to the production model \emph{hurt} ($+0.0014$).\\[3pt]
Q7 & Does draw matter only conditionally & The draw effect on performance is large and monotone (0.25 sd inside to outside). The \emph{market residual} is non-monotone and within 0.2 pp of zero, and a date-atomic draw-bias estimate predicts market error at $r=0.007$. \textbf{Fully priced.}\\[3pt]
Q8 & What the unexplained error looks like & Left-skewed ($-0.729$), mildly heavy-tailed ($+0.810$ excess kurtosis), centre more concentrated than Normal (IQR/sd 1.12 vs 1.35). A two-component mixture wins on held-out likelihood ($+0.056$/obs); Student-$t$ $\nu=5.18$. The bad tail is 4.8$\times$ more likely to carry a severe stewards' phrase.\\[3pt]
Q9 & Is the residual variance constant & \textbf{No.} The largest driver is the market's own price: favourites are more consistent (decile sd ratio 0.784). Joint model oos $R^2$ 0.011 on $\log|e|$, implied p90/p10 spread ratio 1.385.\\[3pt]
Q10 & Are some horses genuinely more inconsistent & \textbf{Yes.} 61.8\% of the spread in per-horse log sd is real; $\tau=0.357$, so 1 sd more volatile means 1.43$\times$ the spread. Split-half persistence $r$ rises to $+0.223$ at $z=7.1$. Forecastable: prior volatility predicts next $|e|$ at $\rho$ up to 0.122.\\[3pt]
Q11 & Race-level shared shocks & \textbf{Large and partly measurable.} 32.1\% of absolute residual variance, $\sim$0.50\,s at 1200\,m. Track speed is a day-level property (within-card $r=0.504$). Pace loads differently by style, correctly signed, $|z|\approx 8$. Market error swings 0.040 on realised pace and 0.015 on projected pace (37.8\% survives).\\[3pt]
Q12 & Reducible against irreducible uncertainty & Clean and monotone. The share attributable to not knowing the horse falls from 19.1\% at 0--3 starts to 2.8\% at 30$+$; the idiosyncratic component is flat at 0.84--0.87 throughout. Beyond $\sim$15 starts nearly everything left is irreducible.\\[3pt]
Q12B & Why prior latent attempts produced little & \textbf{Supported:} weak signal relative to the market; the market already pricing it; final time masking nothing that helps; and --- the new one --- \emph{redundancy with the incumbent's existing columns} (Spearman $+0.39$ between novelty and unhelpfulness). \textbf{Rejected:} poor target choice, insufficient shrinkage, wrong residual distribution, and absence of horse $\times$ environment structure --- all were tested here and none rescues the result.\\[3pt]
Q13 & Can latent distributions give calibrated probabilities & \textbf{Yes, calibrated; no, not competitive.} The simplest variant is better calibrated than the market (RMSE 0.0040 vs 0.0052) and 0.37 of log loss worse. And once the utility scale is fitted, \emph{no} elaboration of the shock structure improves on homoskedastic normal.\\[3pt]
Q14 & Does it add information beyond the market & On one dimension clearly yes --- ability uncertainty predicts the outcome given the price at $t=15$ while being uncorrelated with it. But the incumbent already reconstructs 73\% of it, so the information is in the model even though it is not in the price.\\
\bottomrule
\end{longtable}

\newpage
%======================================================================
\section{Limitations}\label{sec:limitations}
%======================================================================

\begin{itemize}\itemsep3pt
\item \textbf{Prices are retrospective.} Every market probability here is a
result-page price recorded after the pool closed. No bet could have been placed at
it. No dividend was loaded and no return is computed anywhere in this project ---
deliberately, because at this effect size ROI ranks noise.

\item \textbf{No window is unseen.} Every season in the archive has been
development data, in this study and in the reference project. The chronological
splits are real --- a scored year is genuinely after its training rows --- but a
result on 2020--2026 is probability progress, not a locked out-of-sample finding.

\item \textbf{The performance panel is conditioned on finishing.} The sectional
archive records no times for a non-finisher, so 0.60\% of races contribute one
fewer runner here than in the reference table. Every performance distribution in
this document is therefore conditional on completing the race, and the left tail
is truncated at exactly the point where it would be most extreme.

\item \textbf{The within-race target has a mechanical constraint.} $y^{\text{rel}}_z$
is standardised within its own field, so the runners in a race are not independent
by construction and a horse in a weak field scores well for a reason that is not
its own ability. Opponent adjustment --- which the reference project's rank-state
engine attempts --- is not implemented here.

\item \textbf{Jockey and trainer effects are observational.} Assignment is highly
non-random. The within-horse-switch estimate controls for the horse but not for
\emph{why} a jockey was replaced, and the 39\% of apparent jockey variance removed
by horse control is a lower bound on the assignment channel, not an estimate of it.

\item \textbf{Trip effects are largely unobserved.} The stewards' text identifies
severe trouble on 46\% of the worst residuals, which means it misses the rest. A
run compromised in a way nobody wrote down is indistinguishable here from a horse
that ran badly.

\item \textbf{Some contrasts are unanswerable rather than answered.} Horse
$\times$ class sensitivity has too few horses with starts in both class 1--2 and
class 4--5 in both career halves; it is reported as unanswerable. Horse $\times$
field-size sensitivity rests on 36 horses and should not be read at all.

\item \textbf{Non-stationarity is real and only partly handled.} Ability drifts
with a two-month half-life; the stewards' phrasing has changed over nineteen
seasons, which affects the trouble analysis; and the going and course
vocabularies are not stable across the archive.

\item \textbf{Measurement error in the environment propagates.} Going is a
subjective official assessment on an ordinal scale that interleaves turf and
all-weather descriptors, and the par model treats it as a categorical level with
whatever consistency the officials applied.

\item \textbf{The feature test is a development profile, not a canonical run.}
Five years, three seeds, 150 rounds. The reference project's canonical
confirmation is the full chronology at five seeds and 400 rounds; a candidate
whose development standard error is the size of its mean could move either way
there.

\item \textbf{One deviation from the reference project's leakage rule is
reported.} The within-card track-speed signal uses earlier races on the same
meeting. It is forward in time and is not a horse seeing its own result, but it
does breach the letter of a rule the reference project applies uniformly, and it
is flagged wherever it appears rather than folded into a candidate feature.
\end{itemize}

\newpage
%======================================================================
\section{Recommendations}
%======================================================================

\subsection*{Strong evidence --- pursue}

\begin{enumerate}\itemsep3pt
\item \textbf{Recency-weight the horse-ability history at a $\sim$60-day
half-life.} The forward correlation profile is monotone and unambiguous, the
static estimate is the worst of eight tested, and this is already the most direct
open lead in the reference project's own backlog. The measurement here says
\emph{weight, do not truncate}, and gives the weight.

\item \textbf{Build the within-card track-speed estimate and test it properly.}
$r=0.504$ with a race's own shock, rising to 0.584 by the eighth race of a
meeting, is the largest single unexploited signal found. It requires a deliberate
decision about the point-in-time rule --- see the limitation above --- and that
decision should be made explicitly rather than by default. Note it primarily
affects \emph{absolute} time, so its value for win probability must be established
through an interaction, most plausibly with running style.

\item \textbf{Take the sectional finishing band to canonical confirmation.}
$-0.00087$ on the development profile is 20\% of the production edge, with the
right sign in three of five years. Its standard error is the size of its mean, so
the development profile cannot settle it and only a five-seed, 400-round, full-chronology
run can.
\end{enumerate}

\subsection*{Moderate evidence --- worth one more experiment}

\begin{enumerate}\setcounter{enumi}{3}\itemsep3pt
\item \textbf{Projected pace $\times$ running style.} The interaction is real and
correctly signed at $|z|\approx8$, the pace is forecastable at $r=0.41$, and
37.8\% of the market's conditional mispricing survives projection. It did not help
the model ($-0.00028$, 2 of 5 years). The one experiment worth running is a better
pace projection --- the jump from $r=0.03$ to $r=0.41$ on changing the construction
suggests the ceiling is not reached.

\item \textbf{Student-$t$ or mixture shocks in the choice link.} The residual
distribution clearly is not Gaussian ($+0.056$ held-out log-likelihood per
observation for a mixture, $\nu=5.18$ for a $t$). It did not help the generative
model once the scale was fitted, but the reference project's probit link already
gains from swapping Gumbel for normal, and swapping normal for $t$ inside the
\emph{anchored} link is a different and cheaper experiment than the one run here.
\end{enumerate}

\subsection*{Weak or no evidence --- deprioritise}

\begin{enumerate}\setcounter{enumi}{5}\itemsep3pt
\item \textbf{Horse-specific volatility as a mean-model feature.} Real, persistent,
forecastable --- and worth $-0.00011$ with 2 of 5 years improving, and negative in
the generative model. If it has a use it is in \emph{sizing} rather than in the
probability, and staking is out of scope here.

\item \textbf{Horse $\times$ going sensitivity.} Does not persist on the largest
sample of any contrast tested. Do not build it.

\item \textbf{Horse $\times$ venue and $\times$ surface sensitivity as features.}
Genuinely real ($z=8.3$, $z=6.4$) and the block \emph{hurt} the model by $+0.0014$,
the worst of the nine. The gap between real and useful is the finding; a second
construction might do better, but the prior should be low.

\item \textbf{Any draw-bias feature.} Large in performance, fully priced,
and the reference project has already paid $-21$pp for encoding it once.

\item \textbf{Beaten-margin state families.} \texttt{LBW} is the time restated in
quarter-lengths. The eight-feature \texttt{margin\_rank} family and the
\texttt{dynamic\_speed} family measure the same physical quantity, and this is
worth auditing in the reference project on its own terms.
\end{enumerate}

\subsection*{Unknown --- insufficient data}

\begin{enumerate}\setcounter{enumi}{10}\itemsep3pt
\item \textbf{Horse $\times$ jockey compatibility.} Not that it is absent; that
13{,}309 races cannot resolve it. A design with the power to see a true pair sd of
0.10 would need several times the repeated-pairing density this archive has.
\item \textbf{Horse $\times$ class sensitivity.} Too few horses race in both
class 1--2 and class 4--5 within a measurable career.
\end{enumerate}

\subsection*{The shortlist, specified}

\begin{table}[htbp]\centering\small
\caption{Concrete integration shortlist. Deltas are the measured development-profile
paired racewise log-loss change; the production model's whole edge over the market
is $-0.0045$.}
\begin{tabular}{@{}p{0.20\textwidth}p{0.31\textwidth}rp{0.28\textwidth}@{}}
\toprule
Feature & Construction & Measured $\Delta$ & Cost, cadence, safeguard\\
\midrule
\texttt{ability\_recent} &
Date-atomic decayed mean of within-race $z$ of log speed, half-life 60 days,
shrunk by $n/(n+3)$ &
$-0.00102$ &
One pass over the panel; rebuild with the feature table. Read-before-write within
each date makes it date-atomic; no market column enters.\\[4pt]
\texttt{sect\_close} &
Date-atomic prior mean of the final-400\,m section speed residual, expressed
relative to the race's own mean at that band &
$-0.00087$ &
Requires the backward-aligned sectional par (one forward ridge per year). Same
date-atomic guarantee. Needs canonical confirmation.\\[4pt]
Full sectional profile &
The above plus opening, middle, pace balance and profile consistency &
$-0.00118$ &
Five columns; best measured block but its own components disagree, so prefer the
single band until confirmed.\\[4pt]
\texttt{track\_speed\_card} &
Running mean of the par residual over earlier races on the same meeting &
not tested &
Cheap. \textbf{Breaches the date-atomic rule as written} --- requires an explicit
decision, and must be tested through a style interaction since it is common to the
field.\\
\bottomrule
\end{tabular}
\end{table}

\subsection*{A note on how to read a negative result}

By the plan's own definition of success this study succeeds whether or not
anything is promoted, and in fact nothing is. The structural questions have
clear, supported answers; four of them reverse a prior expectation; two prior
results were corrected here after their mechanisms were traced (the pair-effect
anticorrelation artifact, and the unfitted utility scale). The most valuable
output is not a feature but a constraint: at an effect size of $-0.0045$, against
a price this good, and with a model that already reconstructs most of what a new
column would carry, the return on additional orthogonal information is close to
zero. That is worth knowing before building the four hundred and first candidate
column.

\end{document}
